\chapter[Sistema supervisor inteligente]{Sistema supervisor inteligente para optimización del PID local\textsuperscript{\normalfont *}}
\label{ch:supervisor}
\begingroup\renewcommand\thefootnote{*}\renewcommand\theHfootnote{guide.supervisor}\footnotetext{Este capítulo responde al apartado 4.2.d de la guía de la actividad.}\endgroup

El supervisor ocupa el nivel~2 de la jerarquía de control: observa el estado global de la celda y actúa sobre el controlador primario cuando la ocupación de las esteras o la señal de perturbación indica que el régimen nominal ya no es suficiente.

El supervisor es un sistema basado en reglas, organizado como una máquina de estados sobre el nivel de congestión, con tres regímenes (\emph{Normal}, \emph{Riesgo} y \emph{Alarma}): programa las ganancias del PID (\emph{gain scheduling}), ajusta la admisión de piezas y, con un bloque difuso de singletons, decide el reparto del robot entre esteras. Toda su lógica son umbrales y factores fijos, explícitos y auditables, condición necesaria en una capa de seguridad. Por dentro tiene dos niveles: uno de ajuste adaptativo, que modifica las ganancias del PID o condiciona la salida de Fuzzy/DL/QL, y otro de seguridad de capacidad, una barrera dura que actúa siempre para los cuatro controladores.

\section{Variables de supervisión}

El supervisor monitoriza continuamente las siguientes señales en cada muestra $k$:

\begin{longtable}{@{}P{0.20\linewidth}P{0.15\linewidth}P{0.53\linewidth}@{}}
\caption{Variables de supervisión y sus umbrales de actuación.}
\label{tab:sv_variables}\\
\toprule
Variable & Umbral & Descripción y acción del supervisor \\ \midrule \endfirsthead
\toprule Variable & Umbral & Descripción \\ \midrule \endhead
$n_{max}(k) = \max(n_1,n_2)$ & 2.4 (riesgo) & Ajuste suave: $u_i\!+\!0.07$ y admisión $a=0.90$.\\
$n_{max}(k)$ & 2.7 (alarma) & Emergencia de capacidad: $u_i\!+\!0.14$ y admisión $a=0.72$.\\
$d(k) \in \{0,1\}$ & $d > 0$ (sin alarma) & Perturbación activa: $u_i\!+\!0.14$ y admisión $a=0.86$.\\
$b(k) = n_1 - n_2$ & $|b| > 0.15$ & Desequilibrio entre esteras: ajusta reparto $r_1$.\\
$a_{acum}(k)$ & Cualquier & Admisión acumulada; evita bloqueo prolongado.\\
\bottomrule
\end{longtable}

\section{Criterios de ajuste automático del PID}

El supervisor adapta las ganancias del PID en función del estado de congestión de la celda. Se definen tres regímenes de operación:

\begin{longtable}{@{}P{0.13\linewidth}P{0.22\linewidth}P{0.07\linewidth}P{0.07\linewidth}P{0.31\linewidth}@{}}
\caption{Ajuste dinámico de ganancias PID según régimen de operación.}
\label{tab:sv_pid_gains}\\
\toprule
Régimen & Condición & $f_{Kp}$ & $f_{Ki}$ & Justificación \\ \midrule \endfirsthead
\toprule Régimen & Condición & $f_{Kp}$ & $f_{Ki}$ & Justificación \\ \midrule \endhead
Normal & $n_{max}\!\leq\!2.4$ y $d=0$ & $1.0$ & $1.0$ & Control nominal; sin modificación.\\
Riesgo & $2.4<n_{max}\leq2.7$ y $d=0$ & $1.25$ & $0.82$ & Proporcional aumentado para respuesta rápida; integral reducida para evitar windup en la rampa de subida.\\
Alarma & $n_{max}\!>\!2.7$ o $d>0$ & $1.55$ & $0.65$ & Máxima acción proporcional; integral mínima para respuesta de emergencia.\\
\bottomrule
\end{longtable}

Las ganancias ajustadas son $K_p' = K_p \cdot f_{Kp}$ y $K_i' = K_i \cdot f_{Ki}$, con los valores nominales $K_p=1.3$, $K_i=0.18$, $K_d=0.28$ (resintonizados por búsqueda).

\section{Capa de supervisión de capacidad (barrera de seguridad)}

Independientemente del controlador primario activo, la capa de supervisión de capacidad actúa como capa de seguridad que no puede eludirse. Su lógica es:

\begin{align}
\text{Si } n_{max}(k) > N_{alarm}: &\quad u_i(k) = \operatorname{sat}(u_i^* + 0.14,\, 0.05,\, 1.0),\quad a(k) = 0.72 \\
\text{Si } d>0 \text{ (sin alarma de capacidad)}: &\quad u_i(k) = \operatorname{sat}(u_i^* + 0.14,\, 0.05,\, 1.0),\quad a(k) = 0.86 \\
\text{Si } n_{max}(k) > N_{risk}: &\quad u_i(k) = \operatorname{sat}(u_i^* + 0.07,\, 0.05,\, 1.0),\quad a(k) = 0.90 \\
\text{Si no}: &\quad u_i(k) = u_i^*,\quad a(k) = 1.0
\end{align}

El refuerzo de velocidad ($u_i\!+\!0.14$) y el ajuste de ganancias del PID ($f_{Kp}{=}1.55$, $f_{Ki}{=}0.65$)
se activan tanto por alarma de capacidad como por perturbación; solo la fracción de admisión distingue
ambos casos ($a=0.72$ por capacidad crítica frente a $a=0.86$ por perturbación sin alarma).

La Figura~\ref{fig:supervisor_flow} muestra el flujo de decisión del supervisor integrado.

\begin{figure}[H]
\centering
\begin{tikzpicture}[>=Latex, every node/.style={font=\scriptsize}]
  \tikzset{
    start/.style={draw, ellipse, fill=UniPrimary!15, draw=UniPrimaryDeep, minimum width=30mm, minimum height=10mm, align=center},
    dec/.style={draw, diamond, fill=UniLight, draw=UniPrimaryDeep, minimum width=30mm, minimum height=14mm, align=center, aspect=2.4},
    act/.style={draw, rounded corners=1.5mm, fill=UniLight, draw=UniMuted, minimum width=42mm, minimum height=11mm, align=center},
    arr/.style={->, line width=0.85pt, color=UniPrimaryDeep},
    wire/.style={line width=0.85pt, color=UniPrimaryDeep},
    lbl/.style={font=\tiny, fill=white, inner sep=1pt}
  }
  % Columna principal (decisiones y normal)
  \node[start] (in)  at (0, 0)     {Inicio ciclo $k$:\\leer $n_1,n_2,d$};
  \node[dec]   (d1)  at (0,-2.0)   {$n_{max}>2.7$\\o $d>0$?};
  \node[dec]   (d2)  at (0,-4.2)   {$n_{max}>2.4$?};
  \node[act]   (a3)  at (0,-6.3)   {NORMAL: $f_{Kp}{=}1.0$, $f_{Ki}{=}1.0$\\$u$ sin cambio, $a{=}1.0$};
  \coordinate  (mrg) at (0,-7.6);
  \node[dec]   (d3)  at (0,-9.1)   {$|b|>0.15$?};
  \node[start] (out) at (0,-11.1)  {Ejecutar acción\\$u_i,\,r_1,\,a$};
  % Acciones a la derecha
  \node[act]   (a1)  at (5.6,-2.0) {ALARMA/PERT.\\$f_{Kp}{=}1.55$, $f_{Ki}{=}0.65$\\$u{+}{=}0.14$,\; $a{=}0.72/0.86$};
  \node[act]   (a2)  at (5.6,-4.2) {RIESGO\\$f_{Kp}{=}1.25$, $f_{Ki}{=}0.82$\\$u{+}{=}0.07$,\; $a{=}0.90$};
  \node[act]   (a4)  at (5.6,-9.1) {Ajustar reparto:\\$r_1=f(b)$ difuso};
  % Flujo principal
  \draw[arr] (in.south) -- (d1.north);
  \draw[arr] (d1.south) -- node[lbl,right]{NO} (d2.north);
  \draw[arr] (d2.south) -- node[lbl,right]{NO} (a3.north);
  \draw[arr] (a3.south) -- (mrg);
  \draw[arr] (mrg) -- (d3.north);
  \draw[arr] (d3.south) -- node[lbl,right]{NO} (out.north);
  % Ramas SÍ (horizontales, sin cruces)
  \draw[arr] (d1.east) -- node[lbl,above]{SÍ} (a1.west);
  \draw[arr] (d2.east) -- node[lbl,above]{SÍ} (a2.west);
  \draw[arr] (d3.east) -- node[lbl,above]{SÍ} (a4.west);
  % Bus de retorno de a1 y a2 hacia el punto de fusión (mrg)
  \draw[wire] (a1.east) -- (8.0,-2.0);
  \draw[wire] (a2.east) -- (8.0,-4.2);
  \draw[wire] (8.0,-2.0) -- (8.0,-7.6);
  \draw[arr]  (8.0,-7.6) -- (mrg);
  \fill[UniPrimaryDeep] (mrg) circle (1.3pt);
  % Retorno de a4 a la salida (entra por la derecha)
  \draw[arr] (a4.south) |- (out.east);
\end{tikzpicture}
\caption{Flujo de decisión del supervisor inteligente en cada ciclo de muestreo $k$.}
\label{fig:supervisor_flow}
\end{figure}

\simulinkfigure{figures/simulink/sl_supervisor_pid.png}{0.84}{Modelo Simulink del supervisor inteligente que optimiza el PID local (4.2.d): el PID regula la ocupación y el supervisor ajusta sus ganancias por régimen (normal/riesgo/alarma), reduce la admisión y redirige el robot; la capa dura de capacidad permanece siempre activa.}

\simulinkfigure{figures/simulink/sl_supervisor_inteligente.png}{0.84}{Subsistema del supervisor inteligente, encargado de evaluar el estado de ocupación, la perturbación y el desbalance entre esteras para ajustar las ganancias del PID, el caudal efectivo y el reparto del robot.}

\simulinkfigure{figures/simulink/sl_supervisor_pid_ocupacion.png}{0.84}{Subsistema del PID de ocupación, donde se generan las acciones de control $u_1$ y $u_2$ a partir de la referencia $n^*$, los estados $n_1$, $n_2$ y las ganancias ajustadas por el supervisor.}

\simulinkfigure{figures/simulink/sl_supervisor_esteras.png}{0.84}{Subsistema de esteras R2ET, donde se modela la evolución discreta de las ocupaciones $n_1(k)$ y $n_2(k)$ considerando descarga, reparto del robot, caudal efectivo y perturbación.}

\simulinkfigure{figures/simulink/sl_supervisor_scope_ocupacion.png}{0.84}{Respuesta de ocupación del sistema R2ET supervisado: comparación de $n_1(k)$, $n_2(k)$, referencia de 1.5 piezas y capacidad máxima de 3 piezas.}

\simulinkfigure{figures/simulink/sl_supervisor_scope_control.png}{0.84}{Señales de control y supervisión del sistema R2ET: variación del modo de operación, factor de caudal $\alpha$, acciones de control y respuesta ante perturbación.}



\section{Análisis de perturbaciones y fallos}

El supervisor contempla los siguientes modos de perturbación/fallo y sus estrategias de compensación:

\begin{longtable}{@{}P{0.27\linewidth}P{0.25\linewidth}P{0.36\linewidth}@{}}
\caption{Modos de perturbación, síntomas y estrategia de compensación del supervisor.}
\label{tab:sv_faults}\\
\toprule
Perturbación / Fallo & Síntoma & Estrategia del supervisor \\ \midrule \endfirsthead
\toprule Perturbación / Fallo & Síntoma & Estrategia \\ \midrule \endhead
Cambio de peso en piezas (más pesadas) & $n_i$ aumenta rápidamente; $\dot{n}_i > 0$ sostenido. & Activar régimen de alarma: aumentar $u$, reducir admisión $a$, redirigir robot a estera más vacía.\\
Pulso de alta demanda & Ambas $n_i$ suben en paralelo. & Aumentar velocidad de ambas esteras; mantener reparto equilibrado; reducir admisión si $n_{max} > 2.4$.\\
Bloqueo parcial de E1 & $n_1 \uparrow$ mientras $n_2$ estable o baja. & Reducir $r_1$ (menos flujo a E1), aumentar $u_1$ (drenar E1 más rápido).\\
Degradación del actuador & $u_1$ al máximo pero $n_1$ no decrece. & Señal de alarma al SCADA; reducción de $r_1$ a mínimo; activar protocolo de mantenimiento.\\
Pérdida de sensor de $n_i$ & Lectura nula o fuera de rango. & Congelar acción de control del canal afectado; mantener último valor válido; alarma al operador.\\
\bottomrule
\end{longtable}

\section{Coordinación, productividad y adaptación}

\subsection*{Reparto difuso del robot}

El reparto $r_1(k)$ se calcula por promedio ponderado de singletons sobre el desequilibrio $b=n_1-n_2$:

\[
r_1(k) = \operatorname{sat}\!\left(\frac{\mu_{bal} \cdot 0.50 + \mu_{e1h} \cdot 0.30 + \mu_{e2h} \cdot 0.70}{\mu_{bal} + \mu_{e1h} + \mu_{e2h}},\;0.20,\;0.80\right)
\]

Cuando E1 está más cargada ($b>0.15$), $\mu_{e1h}$ domina y la fracción hacia E1 cae a $0.30$, desviando más flujo a E2. Cuando E2 está más cargada, ocurre lo contrario. La saturación garantiza que ambas esteras reciban siempre al menos el $20\%$ del flujo.

\subsection*{Estimación de productividad en tiempo real}

\[
P(k) = \sum_{j=0}^{k}\!\left[\eta\,u_1(j)\,\bigl(1-\ell_1(j)\bigr) + \eta\,u_2(j)\,\bigl(1-\ell_2(j)\bigr)\right]
\]

Si $P(k)$ está por debajo del objetivo programado (por ejemplo, por un atasco prolongado), el supervisor puede relajar la restricción de admisión para recuperar el ritmo de producción tan pronto como la ocupación lo permita.

\subsection*{Q-Learning online}

La tabla Q puede seguir actualizándose durante el despliegue real, con una tasa $\alpha$ baja y exploración ocasional, para adaptarse a deriva lenta del proceso o desgaste de actuadores sin reentrenamiento \emph{offline}. Como muestra la Sección~\ref{sec:ql_online}, esta actualización solo se justifica si existe deriva medible y con una exploración mucho menor ($\varepsilon \ll 0.08$); en operación nominal conviene desplegar la política en modo greedy.
